\documentclass[a4paper]{article}

\input{../etc/preamble.tex}

\title{Appendix A}

\begin{document}

\maketitle

\begin{enumerate} [label=\textbf{\arabic*.}, leftmargin=0em]

      \item Show that the definition of rational equivalence in \SS 1 is equivalent to the equivalence relation generated by the following relation: two cycles $Y, Z$ of codimension $r$ on $X$ are equivalent if there exists a cycle $W$ of codimension $r$ on $X \times \A^1$ which intersects $X \times \{ 0 \}$ and $X \times \{ 1 \}$ properly, and such that $Y = W.(X \times \{ 0 \})$, $Z = W.(X \times \{ 1 \})$.

      \item Prove the following result about Weil divisors, which generalizes (IV, Ex. 2.6), and which is needed to show that $f_*$ is well-defined modulo ratonal equivalence (A3): Let $f : X \to X'$ be a proper, generically finite map of normal varieties, and let $D_1$ and $D_2$ be linearly equivalent Weil divisors on $X$.
      Then $f_* D_1$ and $f_* D_2$ are linearly equivalent Weil divisors on $X'$.

      \item Show directly that any subvariety of degree $d$ in $\P^n$ is rationally equivalent to $d$ times a linear space of the same dimension, by using a projection argument similar to (III, 9.8.3).

      \item Let $\pi : X \to C$ be a ruled surface (V, \SS 2) over a nonsingular curve $C$.
      Show that the group $A^2(X)$ of zero-cycles modulo rational equivalence is isomorphic to $\pic{C}$.

      \item Let $X$ be a surface, let $P \in X$ be a point, and let $\pi : \tilde{X} \to X$ be the monoidal transformation with center $P$ (V, \SS 3).
      Show that $A(\tilde{X}) \cong \pi^* A(X) \oplus \Z$, where $\Z$ is generated by the exceptional curve $E \in A^1(X)$, and the intersection theory is determined by $E^2 = -\pi^*(P)$.

      \item Let $X$ be a nonsingular projective variety of dimension $n$, and let $\Delta \subset X \times X$ be the diagonal.
      Show that $c_n(\fT_X) = \Delta^2$ in $A^n(X)$, under the natural isomorphism of $X$ with $\Delta$.

      \item Let $X$ be a nonsingular projective 3-fold, with Chern classes $c_1, c_2, c_3$.
      Show that
      \begin{equation*}
            1 - p_a = \frac{1}{24} c_1 c_2,
      \end{equation*}
      and for any divisor $D$,
      \begin{equation*}
            \chi(\fO_X(D)) = \frac{1}{12}D.(D - K).(2D - K) + \frac{1}{12} D.c_2 + 1 - p_a.
      \end{equation*}

      \item Let $\fE$ be a locally free sheaf of rank 2 on $\P^3$, with Chern classes $c_1, c_2$.
      Since $A(\P^3) \cong \Z[h]/(h^4)$, we can think of $c_1$ and $c_2$ as integers.
      Show that $c_1 c_2 \equiv 0 \mod{2}$.

      \item \textit{Surfaces in $\P^4$}
      \begin{itemize}
            \item[(a)] Verify the formula of (4.1.3) for the rational cubic scroll in $\P^4$ (V, 2.19.1).
            \item[(b)] If $X$ is a K3 surface in $\P^4$, show that its degree must be 4 or 6.
            \item[(c)] If $X$ is an Abelian surface in $\P^4$, show that its degree must be 10.
            \item[(d)] Determine which of the rational ruled surfaces $X_e$, $e \geq 0$ (V, \SS 2), admit an embedding in $\P^4$.
      \end{itemize}

      \item Use the fact that the tangent sheaf of an Abelian variety is free, to show that it is impossible to embed an Abelian 3-fold in $\P^5$.
\end{enumerate}

\end{document}
